\heading{数据结构与算法A-课程说明}
\noteinfo{以下说明依据本项目现有真题整理，主要反映题库中稳定出现的范围。}

\textbf{一、资料概况}

本项目当前收录本课程题目若干套，并为每套试卷提供 C 语言版与 Python 版。题目来源包括考试回忆、扫描版真题与样题。

\textbf{二、考试范围（按现有题库归纳）}
\begin{itemize}[leftmargin=2em,itemsep=0.2em,topsep=0.2em]
  \item \textbf{线性结构}：顺序表、链表、栈、队列的存储与操作；
  \item \textbf{树与二叉树}：遍历、二叉搜索树、堆、平衡树、Huffman 树；
  \item \textbf{图}：存储（邻接矩阵/表）、遍历（BFS/DFS）、最短路、最小生成树、拓扑排序；
  \item \textbf{查找}：顺序/二分查找、二叉搜索树、散列（哈希）；
  \item \textbf{排序}：插入、选择、冒泡、快速、归并、堆、希尔排序及其稳定性与复杂度；
  \item \textbf{算法设计与分析}：递归、分治、贪心、动态规划、复杂度分析。
\end{itemize}

\textbf{三、试卷结构}
\begin{itemize}[leftmargin=2em,itemsep=0.2em,topsep=0.2em]
  \item \textbf{选择题}：考察基本概念与简单计算，每题 2 分，约 15 题；
  \item \textbf{简答/计算题}：树/图/排序的手工过程，复杂度分析；
  \item \textbf{算法设计/编程题}：伪代码或 C/Python 代码填空与实现。
\end{itemize}

\textbf{四、文件说明}

本目录下源文件命名规则：	exttt{数据结构与算法A-时间-期中/期末/样题-C/P}，其中 C/P 代表使用 C 语言或 Python 语言。扫描版 PDF 为原始试卷；对应的 \texttt{.tex} 文件为转录排版后的版本。C/P 两版除程序代码语言外，题目内容保持一致。

\vfill
\hfill Last updated: \today
